\chapter{Introduction}
\label{chap:intro}

%% -----------------------------------------------------------------------
\section{Multi-Echelon Inventory Optimisation}
\label{sec:intro:problem}

Inventory management decisions lie at the heart of supply-chain efficiency,
resilience, and cost-effectiveness.  Modern supply chains are hierarchical:
goods flow from raw-material suppliers through distribution centres (DCs)
to retail outlets, with inventory held at each stage.  Managing inventory
across these stages — a problem known as \textbf{Multi-Echelon Inventory
Optimisation (MEIO)} — requires making critical replenishment decisions
under demand uncertainty, transportation limitations, and operational
constraints, often without clear visibility into system-wide consequences.

Three compounding difficulties make the MEIO problem harder than its
single-echelon counterpart:

\begin{enumerate}
  \item \textbf{Coupling of stockouts.}  A stockout at a DC propagates
        downstream: retailers awaiting replenishment accumulate unfilled
        demand (backlog), incurring shortage costs.  Simultaneously, the
        DC's depleted inventory triggers a large replenishment order
        upstream, stressing the supplier's capacity.

  \item \textbf{Lead-time amplification.}  The effective replenishment
        lead time seen by a retailer is the cumulative sum of transit
        times across all upstream echelons.  With a 2-day
        supplier-to-DC lead time and a 1-day DC-to-retailer lead time,
        the retailer's effective lead time is 3 days — and any safety
        stock calculation must account for demand variability over the
        entire 3-day window, not just the final mile.

  \item \textbf{The bullwhip effect.}  Small fluctuations in
        end-customer demand are amplified as orders propagate upstream.
        A retailer reacting to a demand spike places a large order at
        the DC; the DC, uncertain of future demand, orders even more from
        the supplier.  This amplification can cause the supplier to
        swing between feast and famine, eroding efficiency across the
        chain \citep{LeeEtAl1997}.
\end{enumerate}

These interdependencies mean that optimising inventory at each node
independently — the standard single-echelon approach — is suboptimal
or even infeasible.  A system-level view is required.

Compounding this challenge, modifying inventory policies in a live supply
chain carries significant operational risk.  Downstream impacts on service
levels, costs, and system stability are difficult to anticipate in advance:
a change to replenishment parameters at the warehouse level may look locally
optimal yet trigger bullwhip-amplified order swings at the supplier, or
deplete retailer safety stock to dangerous levels several weeks later.  This
makes experimentation on production systems costly and often impractical.

%% -----------------------------------------------------------------------
\section{Why a Digital Twin?}
\label{sec:intro:dt}

The classical solution to the MEIO problem is the Clark-Scarf echelon
base-stock policy \citep{ClarkScarf1960}, which is optimal for serial
networks under a set of idealising assumptions:

\begin{itemize}
  \item Backlog symmetry: unmet demand is fully backordered at a known
        linear penalty (no lost sales, no emergency orders).
  \item Infinite supply at the upstream end: the supplier can always
        meet any order placed.
  \item Deterministic lead times: every shipment arrives in exactly
        $L$ days.
  \item No transport capacity: vehicles can carry any quantity at
        the same per-unit cost.
\end{itemize}

Real supply chains routinely violate all four assumptions.  Suppliers have
capacity limits; lead times vary with congestion; transport costs are
non-linear functions of vehicle utilisation.  When transport cost dominates
the budget — as in the system studied in this DDP, where it accounts for
74\% of total cost at the analytical baseline — ignoring vehicle capacity
and consolidation leaves enormous savings on the table.

In practice, decision-makers frequently rely on simulation tools to evaluate
inventory policies and operational scenarios before deployment in order to
avoid the operational risks described above.  However, many available
simulation tools are proprietary, domain-specific, or tightly coupled to
particular industrial settings, limiting their transparency, extensibility,
and suitability for systematic research (Chapter~\ref{chap:litreview} surveys
the landscape in detail).  An open, configurable, MEIO-specific simulation
framework that handles multi-SKU demand, arbitrary network topologies, and
vehicle-capacity-aware transport has not previously been available.

A \textbf{digital-twin simulation} is a computational replica of a supply
chain that can evaluate \emph{any} policy under \emph{any} set of
assumptions.  Unlike closed-form models, it does not need restrictive
assumptions to be tractable.  Unlike real-world pilots, it can evaluate
thousands of policy configurations in hours rather than months, and unlike
proprietary tools, it can be inspected, extended, and adapted to new settings.
The simulation-optimisation paradigm — using a fast simulator as the inner
loop of a black-box optimiser — has emerged as the practical approach for
supply-chain design when analytical solutions are unavailable or
inapplicable \citep{FuEtAl2005}.

%% -----------------------------------------------------------------------
\section{Contributions}
\label{sec:intro:contributions}

This DDP makes the following research contributions:

\begin{enumerate}
  \item \textbf{A modular, tested digital-twin simulation framework for
        MEIO.}  The framework implements seven inventory control policies
        (base-stock, $(s,S)$, $(R,Q)$, $(R,S)$, order-up-to with phase
        offset, $(k,m)$-cycle, and echelon base-stock), an arbitrary
        DAG-topology network model, and a transport layer that enforces
        simultaneous volume and weight capacity constraints with tiered
        cost schedules.  The codebase is validated by approximately 80
        pytest tests covering invariants, policy behaviour, and transport
        mechanics.

  \item \textbf{An Optuna-TPE-based joint optimiser with fill-rate
        constraints.}  The optimiser searches over a mixed-integer space
        of up to 20 parameters (15 base-stock levels + 5 dispatch
        thresholds) and enforces a minimum fill-rate target via a soft
        penalty.  Three modes are supported: inventory-only,
        transport-only, and joint.

  \item \textbf{Empirical demonstration that joint optimisation
        outperforms single-axis optimisation.}  On a 5-SKU, 1N3 network
        with M5 Walmart demand, joint optimisation achieves a 39\% total
        cost reduction relative to the analytical newsvendor baseline while
        simultaneously raising fill rate from 86\% to 95\% — results
        that neither inventory-only nor transport-only optimisation can
        match.
\end{enumerate}

%% -----------------------------------------------------------------------
\section{Thesis Structure}
\label{sec:intro:structure}

The remainder of this thesis is organised as follows.

\textbf{Chapter~\ref{chap:litreview}} reviews the theoretical foundations:
single-echelon newsvendor theory, Clark-Scarf multi-echelon optimality,
periodic-review policies, the bullwhip effect, transport consolidation
models, and simulation-optimisation methods.

\textbf{Chapter~\ref{chap:architecture}} describes the digital-twin
framework architecture: the DAG network model, the discrete-event
simulation loop, the node state machine, the transport layer, the cost
model, and the JSON and CSV configuration pipelines.

\textbf{Chapter~\ref{chap:policies}} presents the seven inventory policies,
their mathematical formulations, and a comparative discussion of their
strengths and limitations.

\textbf{Chapter~\ref{chap:optimization}} describes the Optuna-TPE
optimisation framework: search-space encoding, the fill-rate constraint
mechanism, and the re-evaluation discipline.

\textbf{Chapter~\ref{chap:methodology}} details the experimental design,
dataset (M5 Walmart), warm-up exclusion procedure, multi-seed statistical
validation, and the six main optimisation experiments (E0--E4).

\textbf{Chapter~\ref{chap:results}} presents the results of all
experiments with full cost decompositions, fill-rate breakdowns,
bullwhip ratios, and a cross-experiment comparison.

\textbf{Chapter~\ref{chap:discussion}} synthesises the results, answers
the research questions, draws practical recommendations, and acknowledges
limitations.

\textbf{Chapter~\ref{chap:conclusion}} summarises the contributions,
quantifies the headline results, and outlines the future research roadmap.

%% -----------------------------------------------------------------------
\section{Development Trajectory}
\label{sec:intro:trajectory}

The DDP unfolded in three phases over approximately 1.5 years.

\textbf{Phase 1 — Prototype and validation.}
A single-SKU simulator was built on a simplified 2-warehouse
$\times$ 3-retailer topology.  Four controlled validation experiments
(EV1--EV4) using synthetic and M5 Walmart demand confirmed that the
discrete-event loop, shipment pipeline, and backlog mechanics behave
as predicted by theory.  This phase established confidence in the
simulator's correctness before adding complexity.

\textbf{Phase 2 — Transport layer and scale-up.}
The most significant architectural addition in this phase was the
\textbf{transport module}: a vehicle-capacity-aware layer implementing
dual volume and weight constraints, tiered cost schedules (full-,
half-, and quarter-load rates), and a configurable minimum dispatch
utilisation threshold with a maximum-wait timeout.  This module
transformed the simulator from a pure inventory tool into a
joint inventory-transport digital twin — the distinguishing feature
of the framework.  Concurrently, the network was extended to the
5-SKU 1N3 topology, the inventory policy suite grew from two
implementations (base-stock, $(s,S)$) to seven, and a comprehensive
test suite, multi-seed validation, and bullwhip metric were
instrumented.

\textbf{Phase 3 — Joint optimisation.}
The Optuna-TPE optimiser was integrated and the main experiments
(E0--E4) were designed and executed.  The joint optimiser searches
inventory base-stock levels and transport dispatch thresholds
simultaneously — a search structure only made meaningful by the
transport layer built in Phase~2.  Joint optimisation of inventory
and transport parameters emerged as the central research contribution,
achieving a 39\% cost reduction over the analytical baseline while
exceeding the 92\% fill-rate target.
